\documentclass[12pt]{article}
\usepackage[utf8]{inputenc}
\usepackage[T1]{fontenc}
\usepackage[brazil]{babel}
\usepackage[a4paper,margin=2.5cm]{geometry}
\usepackage{enumitem}
\setlist[itemize]{noitemsep,topsep=4pt}
\title{Sprint 3 - S3\_16\\Front-End Moderno}
\author{Matheus Moreli Durão Silva\\Matrícula 12621CBS217\\Cibersegurança - Bacharelado - Noturno}
\date{2026}
\begin{document}
\maketitle

\section*{Entregável}
Interface Web responsiva com dados dinâmicos. Este documento registra o objetivo do card, a solução desenvolvida e a forma de conferência. Os arquivos de código e as evidências citadas permanecem na mesma pasta deste documento.

\section*{O que foi estudado e aplicado}
\begin{itemize}
\item CSS Grid
\item Media queries
\item Variáveis CSS
\item Arrays e objetos JavaScript
\end{itemize}

\section*{Como conferir}
Sirva a pasta por HTTP e abra \texttt{index.html}. O arquivo \texttt{app.js} percorre um array de objetos e cria três cartões de ativos. O arquivo \texttt{style.css} usa CSS Grid, variáveis CSS e media queries para adaptar o layout. A verificação de execução está em \texttt{evidencia\_execucao.md}.

\section*{Conclusão}
O artefato foi organizado para demonstrar tanto o conceito estudado quanto sua aplicação prática. As evidências de execução só são consideradas válidas quando produzidas pela ferramenta correspondente.
\end{document}
